% Wersja źródłowa: audyt II, jednostka 21
\setcounter{section}{20}
\section{Splot dystrybucji i regularyzacja}
\label{sec:splot21}

W całym rozdziale $d\ge1$, $\K\in\{\R,\C\}$, a dystrybucje
są ciągłymi $\K$-liniowymi funkcjonałami na
$\mathcal D(\R^d;\K)$. Rozpatrujemy splot na całej
$\R^d$; dla ogólnego zbioru otwartego $\Omega$ nie ma
jednej operacji splotu dwóch dystrybucji z
$\mathcal D'(\Omega)$ bez dodatkowych ustaleń.
Oznaczymy przez $\mathcal E'(\R^d)$ przestrzeń dystrybucji
o zwartym nośniku. Nie należy mylić jej ze zbiorem funkcji
testowych $\mathcal D(\R^d)$.

\subsection{Dystrybucje o zwartym nośniku i funkcje gładkie}

\begin{lemat}[Przedłużenie dystrybucji o zwartym nośniku]
\label{lem:compact-action21}
Niech $T\in\mathcal D'(\R^d;\K)$ i
$K=\supp T\Subset\R^d$. Wówczas dla każdej funkcji
$f\in C^\infty(\R^d;\K)$ istnieje dobrze określona wartość
$\langle T,f\rangle:=\langle T,\chi f\rangle$, gdzie
$\chi\in\mathcal D(\R^d)$ jest równa $1$ na otoczeniu $K$.
Wartość ta nie zależy od wyboru $\chi$. Istnieje także
zbiór zwarty $L$ i liczba $m\in\N_0$ oraz $C>0$ takie, że
\[
 |\langle T,f\rangle|\le
 C\max_{|\alpha|\le m}\sup_{x\in L}|\partial^\alpha f(x)|
 \quad(f\in C^\infty(\R^d)).
\]
W szczególności $T$ wyznacza ciągły funkcjonał na
przestrzeni Frécheta $C^\infty(\R^d)$.
\end{lemat}
\begin{proof}
Funkcja $\chi$ istnieje na mocy lematu o gładkim odcięciu.
Jeżeli $\chi_1,\chi_2$ są równe $1$ w otoczeniach $K$,
to $(\chi_1-\chi_2)f$ zeruje się na otoczeniu $K$;
jej nośnik jest rozłączny z $K$, więc z Lematu~\ref{lem:distribution-local19}
wynika $\langle T,(\chi_1-\chi_2)f\rangle=0$.
Gdy $f\in\mathcal D(\R^d)$, identyczny argument dla
$(1-\chi)f$ pokazuje, że przedłużenie zgadza się z pierwotnym
działaniem $T$. Liniowość wynika z użycia jednej ustalonej $\chi$.
Dla ustalonej $\chi$ wybieramy $L=\supp\chi$.
Ciągłość $T$ na $\mathcal D_L$ daje nierówność
$|\langle T,\psi\rangle|\le C_0p_m(\psi)$ dla funkcji
$\psi\in\mathcal D_L$. Podstawiamy $\psi=\chi f$
i stosujemy wzór Leibniza. Dokładniej, dla $|\alpha|\le m$ mamy
$\partial^\alpha(\chi f)=\sum_{\beta\le\alpha}
\binom\alpha\beta(\partial^\beta\chi)(\partial^{\alpha-\beta}f)$.
Poza $L$ wszystkie składniki znikają. Na $L$ ich moduły
ograniczamy przez iloczyny odpowiednich supremów.
Możemy zatem przyjąć
\[
 C=C_0\max\left\{1,\max_{|\alpha|\le m}
 \sum_{\beta\le\alpha}\binom\alpha\beta
 \sup_{x\in\R^d}|\partial^\beta\chi(x)|\right\}.
\]
Otrzymujemy żądane oszacowanie, a zatem ciągłość na
$C^\infty(\R^d)$ także w przypadku $T=0$.
\end{proof}

\begin{uwaga}[Przestrzeń dualna przestrzeni $C^\infty$]
\label{rem:Eprime21}
Odwrotnie, niech $F$ będzie ciągłym funkcjonałem liniowym na
$C^\infty(\R^d;\K)$. Ciągłość daje oszacowanie
$|F(f)|\le C\max_{|\alpha|\le m}\sup_{x\in L}|\partial^\alpha f(x)|$
dla pewnego zwartego $L\subseteq\R^d$, $m\in\N_0$ i $C>0$:
wystarczy zastąpić skończoną rodzinę seminorm jedną seminormą,
biorąc sumę odpowiednich zbiorów zwartych i maksimum rzędów.
Na każdym etapie $\mathcal D_M$ prawa strona jest ograniczona
przez $Cp_m(f)$. Własność uniwersalna granicy induktywnej
(Twierdzenie~\ref{thm:LF-universal18}) daje ciągłość ograniczenia
$T=F|_{\mathcal D}$, a znikanie $F$ na funkcjach testowych
o nośniku rozłącznym z $L$ daje $\supp T\subseteq L$.
Wybierając $\chi=1$ w otoczeniu $L$, mamy
$F(f-\chi f)=0$ z oszacowania, zatem $F(f)=\langle T,\chi f\rangle$.
To pokazuje również jednoznaczność ciągłego przedłużenia.
Otrzymujemy naturalną identyfikację
$\mathcal E'(\R^d)=C^\infty(\R^d)'$, przy czym poniżej
nie będziemy utożsamiać bez wyjaśnienia różnych możliwych
topologii na tej przestrzeni dualnej.
\end{uwaga}

\subsection{Iterowane działania dystrybucji i definicja splotu}

\begin{lemat}[Przestawianie dwóch działań dystrybucji]
\label{lem:iterated-distributions21}
Niech $k,l\ge1$ będą liczbami całkowitymi,
$T\in\mathcal D'(\R^k)$,
$S\in\mathcal D'(\R^l)$ i
$\Phi\in\mathcal D(\R^{k+l})$. Funkcje
\[
 x\longmapsto\langle S_y,\Phi(x,y)\rangle,
 \qquad
 y\longmapsto\langle T_x,\Phi(x,y)\rangle
\]
należą odpowiednio do $\mathcal D(\R^k)$ i
$\mathcal D(\R^l)$. Ponadto
\[
 \langle T_x,\langle S_y,\Phi(x,y)\rangle\rangle
 =\langle S_y,\langle T_x,\Phi(x,y)\rangle\rangle.
\]
\end{lemat}
\begin{proof}
Ustalmy zwarte prostokąty $K_x\subseteq\R^k$ i
$K_y\subseteq\R^l$, których wnętrza zawierają odpowiednie
rzuty $\supp\Phi$. Dla każdego $x$ funkcja $\Phi(x,\cdot)$
należy do $\mathcal D_{K_y}$. Ilorazy różnicowe względem
$x_j$ zbiegają w każdej seminormie $p_m$ tej przestrzeni
do $\partial_{x_j}\Phi(x,\cdot)$: po różniczkowaniu
względem $y$ błąd jest ograniczony przez moduł ciągłości
odpowiedniej pochodnej na iloczynie zwartego otoczenia $x$
i $K_y$, na mocy twierdzenia Newtona--Leibniza.
Powtarzając ten argument dla kolejnych pochodnych i
stosując ciągłość $S$ na $\mathcal D_{K_y}$, otrzymujemy
\[
 \partial_x^\alpha\langle S_y,\Phi(x,y)\rangle
 =\langle S_y,\partial_x^\alpha\Phi(x,y)\rangle.
\]
Ciągłość tych pochodnych wynika z jednostajnej ciągłości
odpowiednich pochodnych $\Phi$ na tych samych zbiorach.
Analogicznie postępujemy z $T$. Nośniki obu otrzymanych
funkcji mieszczą się w odpowiednich rzutach $\supp\Phi$,
więc są to funkcje testowe.

Istnieją $C_T,C_S>0$ i $m_T,m_S\in\N_0$ takie, że
$|\langle T,u\rangle|\le C_Tp_{m_T}(u)$ dla
$u\in\mathcal D_{K_x}$ i analogicznie dla $S$ na
$\mathcal D_{K_y}$. Z powyższego wzoru na pochodne mamy
\[
 |\langle T_x,\langle S_y,\Phi(x,y)\rangle\rangle|
 \le C_TC_S p_{m_T+m_S}(\Phi).
\]
To samo oszacowanie zachodzi dla odwrotnej kolejności.
W szczególności oba działania są ciągłe na każdym etapie
o ustalonym zwartym nośniku.

Pozostaje uzasadnić równość. Dla $\Phi(x,y)=u(x)v(y)$,
gdzie $u,v$ są funkcjami testowymi, obie strony są równe
$\langle T,u\rangle\langle S,v\rangle$.
Pokażemy, że skończone sumy takich funkcji aproksymują
dowolną $\Phi$ w jednym etapie przestrzeni $\mathcal D$.
Wybierzmy nieujemne rzeczywiste funkcje testowe
$\rho$ na $\R^k$ i $\sigma$ na $\R^l$, o całkach $1$
i nośnikach w kulach jednostkowych, przez normalizację
funkcji odcięcia z Lematu~\ref{lem:cutoffs17}.
Niech $\rho_\varepsilon(x)=\varepsilon^{-k}\rho(x/\varepsilon)$
i $\sigma_\varepsilon(y)=\varepsilon^{-l}\sigma(y/\varepsilon)$.
Dla $0<\varepsilon\le1$ określamy
\[
 \Phi_\varepsilon(x,y)=
 \int_{\R^k\times\R^l}\Phi(u,v)
 \rho_\varepsilon(x-u)\sigma_\varepsilon(y-v)
 \,\mathrm d\lambda_{k+l}(u,v).
\]
Wszystkie te funkcje mają nośniki w ustalonym zwartym
prostokącie $L_x\times L_y$, gdzie $L_x,L_y$ zawierają
odpowiednio $K_x+\overline{B(0,1)}$ i
$K_y+\overline{B(0,1)}$. Po zmianie zmiennych całka
jest średnią przesuniętych funkcji $\Phi$ z wagą
$\rho(u)\sigma(v)$. Dla każdego $m$ jednostajna ciągłość
pochodnych $\Phi$ rzędu co najwyżej $m$ daje
$p_m(\Phi_\varepsilon-\Phi)\to0$.

Dla ustalonego $\varepsilon$ całkę definiującą
$\Phi_\varepsilon$ można ograniczyć do $K_x\times K_y$.
Jej sumy Riemanna są skończonymi sumami funkcji postaci
$c\,\rho_\varepsilon(x-u)\sigma_\varepsilon(y-v)$.
Dla każdego skończonego $m$ sumy te zbiegają w $p_m$:
pochodne względem $(x,y)$ można przenieść pod całkę,
a funkcje podcałkowe po różniczkowaniu są jednostajnie
ciągłe na zwartym iloczynie prostokątów zmiennych
$(x,y)$ i $(u,v)$. Poza $L_x\times L_y$ zarówno sumy,
jak i całka znikają. Dla $\varepsilon=1/n$ wybieramy
sumę $F_n$ z błędem $p_n(F_n-\Phi_{1/n})<1/n$.
Wtedy dla każdego ustalonego $m$ i $n\ge m$ mamy
$p_m(F_n-\Phi)\le 1/n+p_m(\Phi_{1/n}-\Phi)\to0$.
Wszystkie $F_n$ mają wspólny zwarty nośnik i są sumami
iloczynów funkcji testowych nad ustalonym ciałem $\K$.
Stosując uzyskane oszacowania do różnicy $F_n-\Phi$,
przechodzimy do granicy w równości obu działań na $F_n$.
\end{proof}

\begin{lemat}[Operator testowy indukowany dystrybucją]
\label{lem:convolution-transpose21}
Niech $T\in\mathcal E'(\R^d)$, $K=\supp T$, i niech
$\varphi\in\mathcal D(\R^d)$. Funkcja
\[
 (A_T\varphi)(y)=\langle T_x,\varphi(x+y)\rangle
\]
należy do $\mathcal D(\R^d)$ i spełnia
$\supp(A_T\varphi)\subseteq\supp\varphi-K$.
Operator $A_T:\mathcal D(\R^d)\to\mathcal D(\R^d)$
jest liniowy i ciągły, a obrazy zbiorów ograniczonych
są ograniczone.
\end{lemat}
\begin{proof}
W działaniu $T_x$ korzystamy z przedłużenia z
Lematu~\ref{lem:compact-action21}. Dla $y$ z ustalonego
zbioru zwartego funkcje $x\mapsto\varphi(x+y)$ mają
wspólny zwarty nośnik; można więc różniczkować po $y$:
$\partial_y^\beta(A_T\varphi)(y)=
\langle T_x,\partial^\beta\varphi(x+y)\rangle$.
Jeżeli $y\notin\supp\varphi-K$, funkcja
$\varphi(\cdot+y)$ zeruje się na otoczeniu $K$,
stąd wartość działania $T$ jest zerowa. Zatem obraz
jest funkcją testową.

Dla ustalonego $L\Subset\R^d$ i
$\varphi\in\mathcal D_L$ oszacowanie z
Lematu~\ref{lem:compact-action21}, zastosowane do
przesunięć funkcji $\varphi$, daje
\[
 p_m(A_T\varphi)\le C_{T,L,m}\,p_{m+r}(\varphi),
 \qquad \supp(A_T\varphi)\subseteq L-K,
\]
gdzie $r$ jest rzędem oszacowania z Lematu~\ref{lem:compact-action21}.
Istotnie, dla $|\beta|\le m$ stosujemy to oszacowanie do
$f(x)=\partial^\beta\varphi(x+y)$; każda pochodna $f$
rzędu co najwyżej $r$ ma moduł ograniczony przez
$p_{m+r}(\varphi)$ niezależnie od $y$.
Można więc wziąć stałą z tego lematu niezależną od $L,m$.
Zbiór $L-K$ jest zwarty jako obraz $L\times K$ przez
ciągłe odwzorowanie $(u,v)\mapsto u-v$; dla $K=\varnothing$
operator jest zerowy.
To dowodzi ciągłości na każdym etapie $\mathcal D_L$,
a z własności uniwersalnej granicy induktywnej
wynika ciągłość na $\mathcal D$. Ciągłe odwzorowanie liniowe
przenosi zbiory ograniczone na ograniczone: dla otoczenia
zera $U$ w przestrzeni docelowej $A_T^{-1}(U)$ jest
otoczeniem zera, a z $B\subseteq t A_T^{-1}(U)$
wynika $A_T(B)\subseteq tU$.
\end{proof}

\begin{definicja}[Splot dystrybucji]
\label{def:convolution21}
Niech $T\in\mathcal E'(\R^d)$ i
$S\in\mathcal D'(\R^d)$. Ich \emph{splotem} nazywamy
dystrybucję $T*S\in\mathcal D'(\R^d)$ określoną wzorem
\[
 \langle T*S,\varphi\rangle
 =\langle S_y,\langle T_x,\varphi(x+y)\rangle\rangle
 =\langle S,A_T\varphi\rangle,
 \qquad\varphi\in\mathcal D(\R^d).
\]
Jeżeli tylko drugi czynnik ma zwarty nośnik, definiujemy
$S*T:=T*S$. Gdy oba czynniki mają zwarty nośnik,
zgodność obu możliwych kolejności zostanie dowiedziona
poniżej. W szczególności nie definiujemy tutaj splotu
dowolnych dwóch dystrybucji bez warunków na nośniki.
\end{definicja}

\begin{twierdzenie}[Istnienie, ciągłość i nośnik splotu]
\label{thm:convolution-properties21}
Niech $T\in\mathcal E'(\R^d)$ i
$S\in\mathcal D'(\R^d)$. Wówczas:
\begin{enumerate}
\item splot $T*S$ jest dobrze określoną dystrybucją;
\item $\supp(T*S)\subseteq\supp T+\supp S$;
\item dla ustalonego $T$ odwzorowanie
$S\mapsto T*S$ jest ciągłym operatorem liniowym na
$\mathcal D'(\R^d)$ zarówno dla topologii słabej-*, jak
i dla topologii silnej;
\item jeżeli także $S$ ma zwarty nośnik, to $T*S$ ma
zwarty nośnik.
\end{enumerate}
\end{twierdzenie}
\begin{proof}
Pierwszy punkt jest konsekwencją ciągłości $A_T$:
złożenie ciągłego funkcjonału $S$ z $A_T$ jest ciągłym
funkcjonałem na $\mathcal D$. Oznaczmy $K=\supp T$, $F=\supp S$. Suma $K+F$
jest domknięta: jeżeli $k_n+f_n\to z$, $k_n\in K$,
$f_n\in F$, to z ciągu $(k_n)$ wybieramy podciąg zbieżny
do $k\in K$, a odpowiedni podciąg $(f_n)$ zbiega do
$z-k\in F$. Dla pustego składnika suma jest pusta.
Jeżeli $\supp\varphi$ jest rozłączny z $K+F$, to
$(\supp\varphi-K)\cap F=\varnothing$, gdyż
$y=u-k\in F$ oznaczałoby $u=k+y\in K+F$.
Z Lematu~\ref{lem:convolution-transpose21} wynika więc
$\supp(A_T\varphi)\cap F=\varnothing$, a z
Lematu~\ref{lem:distribution-local19} otrzymujemy
$\langle S,A_T\varphi\rangle=0$.
Stosując to do wszystkich funkcji testowych o nośniku
w otwartym dopełnieniu $K+F$, dostajemy
$\supp(T*S)\subseteq K+F$.

Dla ustalonej funkcji testowej $\varphi$ zachodzi
$|\langle T*S,\varphi\rangle|=
|\langle S,A_T\varphi\rangle|$, co zapewnia
ciągłość słabą-*. Dla zbioru ograniczonego
$B\subseteq\mathcal D$ mamy
\[
 \sup_{\varphi\in B}|\langle T*S,\varphi\rangle|
 =\sup_{\psi\in A_T(B)}|\langle S,\psi\rangle|.
\]
Lemat~\ref{lem:convolution-transpose21} daje
ograniczoność $A_T(B)$, a więc ciągłość silną.
Jeśli oba nośniki są zwarte, ich suma jest zwarta,
co daje ostatni punkt.
\end{proof}

\subsection{Przemienność, łączność i różniczkowanie}

\begin{twierdzenie}[Przemienność i łączność splotu]
\label{thm:convolution-algebra21}
Niech $T,S,U\in\mathcal D'(\R^d)$.
Jeśli co najmniej jeden z czynników $T,S$ ma zwarty nośnik,
to $T*S=S*T$, a wynik nie zależy od sposobu użycia
Definicji~\ref{def:convolution21}. Jeśli co najmniej dwa
spośród $T,S,U$ mają zwarty nośnik, to
\[
 (T*S)*U=T*(S*U).
\]
W szczególności $\mathcal E'(\R^d)$ jest przemienną
algebrą z jedynką $\delta_0$ względem splotu,
a $\mathcal D'(\R^d)$ jest modułem nad tą algebrą.
\end{twierdzenie}
\begin{proof}
Jeżeli $T$ i $S$ mają zwarty nośnik, wybieramy funkcje
odcięcia $\chi_T,\chi_S$, równe $1$ na otoczeniach
odpowiednich nośników. Dla $\varphi\in\mathcal D$ funkcja
$\Phi(x,y)=\chi_T(x)\chi_S(y)\varphi(x+y)$
należy do $\mathcal D(\R^{2d})$. Oba możliwe
obliczenia splotu są iterowanymi działaniami na $\Phi$;
równość wynika z Lematu~\ref{lem:iterated-distributions21}.
Jeśli tylko jeden czynnik ma zwarty nośnik,
przemienność jest częścią definicji, a powyższy
argument zapewnia jej niesprzeczność.

Dla łączności można, bez utraty ogólności, założyć,
że $T$ i $S$ mają zwarty nośnik. Dla ustalonej
$\varphi\in\mathcal D$ obie strony są trzema
iterowanymi działaniami na $\varphi(x+y+z)$.
Wstawiamy funkcje odcięcia w zmiennych $x$ i $y$
równe $1$ na otoczeniach nośników $T$ i $S$,
a w zmiennej $z$ funkcję odcięcia równą $1$
na zbiorze $\supp\varphi-\supp T-\supp S$.
Uzyskujemy funkcję testową trzech zmiennych.
Kolejne zastosowania Lematu~\ref{lem:iterated-distributions21}
pozwalają dowolnie zmieniać kolejność działań.
To daje łączność; pozostałe ustawienia czynników
wynikają z przemienności. Wzory
$\delta_0*T=T$ oraz $T*(S+U)=T*S+T*U$
wynikają bezpośrednio z definicji.
\end{proof}

\begin{definicja}[Translacja dystrybucji]
\label{def:translation21}
Dla $a\in\R^d$ i $T\in\mathcal D'(\R^d)$ definiujemy
$\tau_aT$ wzorem
$\langle\tau_aT,\varphi\rangle=
\langle T,\varphi(\cdot+a)\rangle$.
Dla dystrybucji regularnej $T_f$ odpowiada to
przesunięciu funkcji $f(\cdot-a)$.
\end{definicja}

\begin{przyklad}[Dystrybucje Diraca i przesunięcia]
\label{ex:delta-convolution21}
Dla $a,b\in\R^d$ i $T\in\mathcal D'$ mamy
$\delta_a*T=\tau_aT$, $\delta_a*\delta_b=\delta_{a+b}$.
Dla wielowskaźników $\alpha,\beta\in\N_0^d$ otrzymujemy
$\partial^\alpha\delta_0*\partial^\beta\delta_0
=\partial^{\alpha+\beta}\delta_0$.
\end{przyklad}

\begin{twierdzenie}[Pochodne splotu]
\label{thm:convolution-derivative21}
Niech $T,S\in\mathcal D'(\R^d)$ i niech co najmniej
jeden z nich ma zwarty nośnik. Dla każdego
$\alpha\in\N_0^d$ zachodzi
\[
 \partial^\alpha(T*S)
 =(\partial^\alpha T)*S
 =T*(\partial^\alpha S).
\]
W szczególności dla wielomianu $P$ w $d$ zmiennych
i operatora stałowspółczynnikowego $P(\partial)$
operator splotu z ustaloną dystrybucją o zwartym
nośniku komutuje z $P(\partial)$.
\end{twierdzenie}
\begin{proof}
Załóżmy, że $T$ ma zwarty nośnik. Dla
$\varphi\in\mathcal D$ i $j\in\{1,\ldots,d\}$ mamy
\[
 \begin{aligned}
 \langle\partial_j(T*S),\varphi\rangle
 &=-\langle S_y,\langle T_x,\partial_j\varphi(x+y)\rangle\rangle\\
 &=\langle S_y,\langle(\partial_jT)_x,\varphi(x+y)\rangle\rangle
 =\langle(\partial_jT)*S,\varphi\rangle.
 \end{aligned}
\]
W drugim kroku różniczkowanie dotyczy zmiennej $x$.
Aby uzyskać $T*(\partial_jS)$, zauważamy, że
$\partial_{y_j}A_T\varphi=A_T(\partial_j\varphi)$,
więc definicja pochodnej $S$ daje tę samą wartość.
Indukcja dowodzi wzoru dla $\alpha$. Gdy zwarty
nośnik ma tylko $S$, stosujemy przemienność.
\end{proof}

\begin{przyklad}[Inkluzja nośników może być ostra]
\label{ex:support-strict21}
Niech $T=\delta_0-\delta_1\in\mathcal E'(\R)$,
a $S=T_1$ będzie dystrybucją odpowiadającą funkcji
stałej $1$. Wtedy $T*S=T_1-T_1=0$,
mimo że $\supp T+\supp S=\R$. Dlatego
inkluzji nośników nie można zastąpić równością.
\end{przyklad}

\subsection{Splot z funkcją testową i aproksymacja jedności}

\begin{twierdzenie}[Wygładzanie dystrybucji]
\label{thm:smoothing21}
Niech $T\in\mathcal D'(\R^d)$ i
$\rho\in\mathcal D(\R^d)$. Wówczas splot
$T*\rho$ jest dystrybucją regularną odpowiadającą
funkcji $u\in C^\infty(\R^d)$ danej przez
\[
 u(x)=\langle T_y,\rho(x-y)\rangle.
\]
Dla każdego wielowskaźnika $\alpha$ zachodzi
$\partial^\alpha u(x)
=\langle T_y,(\partial^\alpha\rho)(x-y)\rangle$.
Ponadto $\supp u\subseteq\supp T+\supp\rho$.
\end{twierdzenie}
\begin{proof}
Przy $x$ należącym do ustalonego zbioru zwartego
rodzina funkcji testowych $y\mapsto\rho(x-y)$
ma wspólny zwarty nośnik. Różniczkowanie po parametrze
$x$ zachodzi więc w topologii $\mathcal D$, skąd
gładkość $u$ i wzory na pochodne. Jeśli $x$
nie należy do $\supp T+\supp\rho$, funkcja
$y\mapsto\rho(x-y)$ znika na otoczeniu $\supp T$,
a zatem $u(x)=0$. Równość dystrybucji $T*\rho=T_u$
sprawdzamy, stosując Lemat~\ref{lem:iterated-distributions21}
do funkcji testowej $\rho(x-y)\varphi(x)$:
\[
 \begin{aligned}
 \int_{\R^d}u(x)\varphi(x)\,\mathrm d\lambda_d(x)
 &=\left\langle T_y,
      \int_{\R^d}\rho(x-y)\varphi(x)\,\mathrm d\lambda_d(x)
   \right\rangle\\
 &=\langle T,\check\rho*\varphi\rangle
 =\langle T*\rho,\varphi\rangle,
 \end{aligned}
\]
gdzie $\check\rho(z)=\rho(-z)$. W pierwszym
kroku zamiana całkowania i działania dystrybucji
jest dozwolona również bez globalnego zwartego nośnika
$T$, ponieważ $\varphi$ i $\rho$ mają zwarte nośniki.
\end{proof}

\begin{definicja}[Standardowa aproksymacja jedności]
\label{def:mollifier21}
Wybierzmy $\rho\in\mathcal D(\R^d)$ taką, że
$\rho\ge0$, $\supp\rho\subseteq\overline{B(0,1)}$
oraz $\int_{\R^d}\rho(x)\,\mathrm d\lambda_d(x)=1$.
Dla $\varepsilon>0$ piszemy
$\rho_\varepsilon(x)=\varepsilon^{-d}\rho(x/\varepsilon)$.
\end{definicja}

\begin{twierdzenie}[Regularyzacja jest silnie zbieżna]
\label{thm:strong-mollification21}
Dla każdego $T\in\mathcal D'(\R^d)$ i każdej
standardowej aproksymacji jedności z
Definicji~\ref{def:mollifier21} dystrybucje
$T*\rho_\varepsilon$ są reprezentowane przez funkcje
$C^\infty(\R^d)$ oraz
\[
 T*\rho_\varepsilon\longrightarrow T
 \quad\text{w topologii silnej }\beta(\mathcal D',\mathcal D)
 \quad(\varepsilon\downarrow0).
\]
Jeśli $T$ ma zwarty nośnik $K$, to
$\supp(T*\rho_\varepsilon)\subseteq K+
\overline{B(0,\varepsilon)}$, więc dla małych
$\varepsilon$ funkcje wygładzające mają nośniki
w jednym zbiorze zwartym.
\end{twierdzenie}
\begin{proof}
Z Twierdzenia~\ref{thm:smoothing21} mamy wygładzanie.
Dla $\varphi\in\mathcal D$ zachodzi
\[
 \langle T*\rho_\varepsilon-T,\varphi\rangle
 =\langle T,\check\rho_\varepsilon*\varphi-\varphi\rangle.
\]
Niech $B\subseteq\mathcal D$ będzie ograniczony.
Z twierdzenia o ograniczoności w $\mathcal D$
wynika istnienie $K\Subset\R^d$ zawierającego nośniki
wszystkich $\varphi\in B$, przy czym
$M_m(B)=\sup_{\varphi\in B}p_m(\varphi)<\infty$
dla wszystkich $m$. Dla $0<\varepsilon\le1$
różnice $\check\rho_\varepsilon*\varphi-\varphi$
mają nośniki w jednym zwartym zbiorze
$L=K+\overline{B(0,1)}$. Twierdzenie o wartości
średniej zastosowane do pochodnych rzędu $\le m$ daje
\[
 \sup_{\varphi\in B}p_m(\check\rho_\varepsilon*\varphi
 -\varphi)
 \le C_{\rho,d}\varepsilon\,M_{m+1}(B),
 \quad
 C_{\rho,d}=\int_{\R^d}\|z\|_1\rho(z)\,\mathrm d\lambda_d(z).
\]
Ciągłość $T|_{\mathcal D_L}$ oznacza, że dla pewnych
$m,C$ mamy $|\langle T,\psi\rangle|\le Cp_m(\psi)$
gdy $\psi\in\mathcal D_L$. Dlatego
\[
 \sup_{\varphi\in B}|\langle T*\rho_\varepsilon-T,
 \varphi\rangle|
 \le CC_{\rho,d}\varepsilon M_{m+1}(B)\longrightarrow0.
\]
Ponieważ $B$ jest dowolnym zbiorem ograniczonym,
jest to zbieżność silna. Oszacowanie nośnika wynika
z Twierdzenia~\ref{thm:smoothing21}.
\end{proof}

\begin{wniosek}[Gęstość funkcji testowych dla zwartych nośników]
\label{cor:dense-test-compact21}
Niech $T\in\mathcal E'(\R^d)$. Istnieje ciąg
$(f_n)\subseteq\mathcal D(\R^d)$ taki, że
$T_{f_n}\to T$ silnie w $\mathcal D'(\R^d)$.
Można wziąć $f_n=T*\rho_{1/n}$.
\end{wniosek}
\begin{proof}
Funkcje $f_n$ są gładkie i mają nośniki zwarte
z poprzedniego twierdzenia. Zbieżność silna także
wynika z tego twierdzenia.
\end{proof}

\begin{przyklad}[Regularyzacja delty Diraca]
\label{ex:delta-smooth21}
Dla $T=\delta_a$ zachodzi
$\delta_a*\rho_\varepsilon=\rho_\varepsilon(\cdot-a)$.
Funkcje te stają się coraz bardziej skupione
w otoczeniu $a$, a jako dystrybucje zbiegają
silnie do $\delta_a$. Nie ma tu zbieżności w
normie $L^1_{\mathrm{loc}}$ do funkcji:
$\delta_a$ nie jest dystrybucją regularną
reprezentowaną przez funkcję lokalnie całkowalną.
\end{przyklad}

\subsection{Ograniczenia operacji splotu i powiązania z dalszą teorią}

\begin{uwaga}[Kiedy splot nie jest zdefiniowany]
\label{rem:convolution-limits21}
Warunek zwartego nośnika przynajmniej jednego
czynnika jest wystarczający, ale niekonieczny.
Dla funkcji odpowiednio całkowalnych splot może istnieć
bez tego warunku. Natomiast dla dowolnych dwóch
dystrybucji nie istnieje jedna kanoniczna operacja
splotu o wszystkich powyższych własnościach.
Na przykład formalny splot dwóch funkcji stałych
na $\R^d$ wymagałby całkowania stałej $1$
po całej $\R^d$, co nie daje skończonej wartości.
Dalsze rozszerzenia opierają się między innymi
na warunkach na nośniki umożliwiających właściwe
rzutowanie ich sum lub na dodatkowych własnościach
analitycznych czynników.
\end{uwaga}

\begin{uwaga}[Dlaczego jest to zagadnienie analizy funkcjonalnej]
\label{rem:convolution-tvs21}
Kluczowe kroki nie polegały na rozwiązywaniu równań
różniczkowych. Użyliśmy ciągłości operatora
$A_T$ na przestrzeni funkcji testowych, transpozycji
na jej przestrzeń dualną, przenoszenia zbiorów ograniczonych
oraz zbieżności silnej w tej przestrzeni dualnej.
Są to konstrukcje z teorii przestrzeni liniowo-topologicznych.
Dalsze zastosowania splotu do rozwiązywania równań
różniczkowych pozostawiamy odrębnemu skryptowi.
\end{uwaga}
